\section{Ramified geometric trace and the original-zeta dual
action}\label{ramified-geometric-trace-and-the-original-zeta-dual-action}

Complete derivation, 24 September 2026. Proof locators GTR0--GTR11.

\subsection{GTR0. Construction stage and
sources}\label{gtr0.-construction-stage-and-sources}

The support remains \(\tau\langle Z_1;\mathrm{no}\ Z_2\rangle\). The
user's complete arithmetic reconstruction precedes every coefficient
space, scalar, positive integer, and analytic coordinate below. No
addition or numerical coordinate is assigned to \(\tau\). The governing
corpus and prerequisite rule and correction chains were read; the full
passages rejecting an added distance/vector and the pooling of separate
branch quotients were reread. The geometric sheets of an already
constructed cover are not pooled arithmetic histories. The original two
labelled support charts remain separate throughout.

The inputs are the complete CSP0--CSP13 sphere construction and
CDW0--CDW10 explicit derived direct image in the receiving task. Their
exact coefficients, restriction maps, periods and positive integration
conventions are retained. The continuous transpose comparison is
CDI0--CDI12, and the original-zeta residue comparison is RD0--RD9. These
proofs were read in full. No continuous-dual complex is identified here
with an ordinary Verdier dual without a separate constructed comparison.

The literature comparison is the Stacks Project,
\href{https://stacks.math.columbia.edu/tag/0GKE}{tag 0GKE}, especially
the trace existence, composition and flat quasi-finite weighting
statements. Original author TeX \texttt{more-etale.tex},
lines1761--2127, was read at pinned commit
\nolinkurl{a04446e57ec1fbc252a871afcec7752fb2807b14}; its source receipt
is retained locally. That source concerns étale sheaves. The topological
sheaf trace used here is constructed directly below, so its existence is
not inferred by changing categories in the cited theorem. The original
Connes--Consani source identity and reading coverage remain CSP0's
\href{https://arxiv.org/abs/0903.2024v3}{arXiv:0903.2024v3}, §5, and the
later signed chart construction. Deligne's comparison remains
\href{https://numdam.org/item/PMIHES_1980__52__137_0/}{Weil II}, §3.6,
printed213--214.

\subsection{GTR1. The actual sheaf and positive-degree
covers}\label{gtr1.-the-actual-sheaf-and-positive-degree-covers}

Use \(Y=\mathbb P^1(\mathbb C)\), \(D=\{0,\infty\}\),
\(Y^\circ=\mathbb C^\times\), and the CSP coefficient spaces \[
S=\{f\in\mathcal S(\mathbb R;\mathbb C):f(-x)=f(x),\ f(0)=0,\ \int_{\mathbb R}f=0\},
\quad V_p=S\oplus E_p,\quad E_p=\mathbb C^2\quad(p=+,-),
\] \[
A=\{a\in C^\infty(\mathbb R_{>0}):
\sup_{u>0}(u^N+u^{-N})|(u\partial_u)^ja(u)|<\infty
\text{ for every }N,j\ge0\}.
\tag{GTR1.1}
\] The full restrictions are \[
\Sigma f(u)=2\sum_{k\ge1}f(ku),\quad
Ra(u)=u^{-1}a(u^{-1}),\quad
r_+(f,e)=\Sigma f,\quad r_-(h,e)=R\Sigma h.
\tag{GTR1.2}
\] Write \(J=\Sigma S\), \(Q=A/J\). The actual source-image theorem
proves closedness of \(J\) and the continuous inverse of
\(\Sigma:S\to J\). With \(i:D\hookrightarrow Y\), \[
\mathscr F=\underline A_Y\times_{i_*(A\oplus A)}i_*(V_+\oplus V_-).
\tag{GTR1.3}
\] On a connected open meeting one pole, a section consists of a
constant \(a\in A\) and \(v_p\in V_p\) with \(a=r_pv_p\); on an interior
connected open it is a constant \(a\). This follows directly from the
fibre product, and specifies the local maps used below.

For every recovered integer \(n\ge1\), let \(b_n(z)=z^n\), with both
poles fixed. The equality \(\pi b_n=\pi\) supplies the canonical
identification \(b_n^{-1}\mathscr F=\mathscr F\). This is an
identification of sheaves pulled from the same three-point diagram, not
an inverse geometric map.

The local multiplicities are \[
e_n(x)=\begin{cases}n,&x=0\text{ or }\infty,\\1,&x\in\mathbb C^\times,\end{cases}
\qquad \sum_{b_n(x)=y}e_n(x)=n.
\tag{GTR1.4}
\] At zero the local ring quotient is \(\mathbb C[z]_{(z)}/(z^n)\), of
length \(n\) with basis \(1,z,\ldots,z^{n-1}\); the coordinate \(w=1/z\)
gives the same calculation at infinity. Away from the poles the
derivative \(nz^{n-1}\) is nonzero, and there are exactly \(n\) distinct
roots. Thus (GTR1.4) includes both ramified points and every interior
sheet. It agrees with the flat quasi-finite weighting in the cited
source, without removing the ramification.

\subsection{GTR2. The weighted trace exists on the actual topological
sheaf}\label{gtr2.-the-weighted-trace-exists-on-the-actual-topological-sheaf}

Define \[
\operatorname{Tr}_n:b_{n*}b_n^{-1}\mathscr F=b_{n*}\mathscr F\longrightarrow\mathscr F
\tag{GTR2.1}
\] as follows. On a sufficiently small interior disk the inverse image
is a disjoint union of \(n\) disks, and the map on sections is \[
(a_1,\ldots,a_n)\longmapsto\sum_{j=1}^n a_j.
\tag{GTR2.2}
\] On a sufficiently small disk about a pole, the inverse image is one
pole disk and a section is \((r_pv_p,v_p)\). Its image is \[
(r_pv_p,v_p)\longmapsto(nr_pv_p,nv_p).
\tag{GTR2.3}
\] The restriction of (GTR2.3) to an interior subdisk equals (GTR2.2):
all \(n\) sheet values are the restriction \(r_pv_p\). Reordering the
sheets does not change the sum. These formulas therefore agree on a
basis of neighborhoods and every overlap, and define a sheaf morphism.
Each local map is a finite continuous linear operation on the specified
Fréchet spaces. At each pole both endpoint coordinates are multiplied by
\(n\); they are not discarded because their generic restriction
vanishes.

These local descriptions determine the map uniquely once its full
weighted stalk values are specified. Generic-sheet agreement alone would
not determine a map on the endpoint kernels \(E_p\). Formula (GTR2.3)
supplies that missing value explicitly from the local multiplicity.

The pullback unit \(u_n:\mathscr F\to b_{n*}b_n^{-1}\mathscr F\) is
diagonal on an interior disk and is the identity on a pole stalk. Hence
\[
\boxed{\operatorname{Tr}_n u_n=n\,1_{\mathscr F}.}
\tag{GTR2.4}
\] This is an equality on every stalk and therefore on sheaves. No
division by \(n\) is used in constructing the trace.

The finite proper map \(b_n\) has exact direct image on vector-space
sheaves. To see this without a finite-rank hypothesis, its direct-image
stalk at \(y\) is the finite direct sum of the stalks at
\(b_n^{-1}(y)\). Indeed small disjoint neighborhoods of the finite fibre
contain the inverse image of some neighborhood of \(y\): otherwise
properness supplies an accumulation point outside those neighborhoods
mapping to \(y\), a contradiction. Taking the corresponding neighborhood
limits proves the stalk formula. Finite direct sums and stalk functors
are exact. This also proves \(Rb_{n*}=b_{n*}\); derived global sections
of the pushforward equal those on the source. Thus (GTR2.1) gives an
actual transfer on all cohomology groups.

The formula is natural in every continuous morphism of coefficient
diagrams, since such maps commute with finite sums and multiplication by
\(n\). In particular it is natural for the actual sequence with overlap
coefficients \(J\to A\to Q\), retaining the entire \(J\). Composition is
exact: grouping roots in the composite cover multiplies their local
multiplicities, \[
e_{mn}(x)=e_n(x)e_m(b_n(x)),
\] and the two successive finite sums equal the single weighted sum.
Consequently the transfer for \(b_{mn}\) equals the composite transfer
for \(b_n\) and \(b_m\), with the direct-image identifications retained.

\subsection{GTR3. Transfer on the full cohomology and support
diagrams}\label{gtr3.-transfer-on-the-full-cohomology-and-support-diagrams}

Let \(P_n\) denote geometric pullback alone. The positive angular form
is \[
\vartheta=\frac{d\arg z}{2\pi},\qquad b_n^*\vartheta=n\vartheta.
\tag{GTR3.1}
\] The positive sphere form and its full integral are \[
\omega_Y=\frac{dx\wedge dy}{\pi(1+|z|^2)^2},\quad
\int_Yb_n^*\omega_Y
=2n^2\int_0^\infty\frac{r^{2n-1}}{(1+r^{2n})^2}\,dr=n.
\tag{GTR3.2}
\] CSP and CDW give the explicit global coefficient complex \[
D^0=V_+\oplus V_-\xrightarrow{r_+-r_-}D^1=A\xrightarrow0D^2=A.
\tag{GTR3.3}
\] Its groups are \(H^0=Z=\ker(r_+-r_-)\), \(H^1=Q\), \(H^2=A\). On this
model geometric pullback is \((1,1,n)\). In degree one it is the
identity because the original degree-zero overlap constants are
unchanged; in degree two (GTR3.2) gives the factor. These are invertible
operators on these already constructed complex coefficient groups.

Taking cohomology in (GTR2.4) and using the proved pullback values
determines the actual transfer: \[
\begin{array}{c|ccc}
&H^0(Y,\mathscr F)=Z&H^1(Y,\mathscr F)=Q&H^2(Y,\mathscr F)=A\\\hline
P_n&1&1&n\\
\operatorname{Tr}_n&n&n&1.
\end{array}
\tag{GTR3.4}
\] The map \((n,n,1)\) on (GTR3.3) is a continuous cochain
representative of these transfer values: its only nonzero differential
commutes with \(n\). No assertion that the map \(z\mapsto z^n\) has a
single-valued inverse is involved. This statement identifies the induced
cohomology maps and an explicit compatible coefficient representative;
it does not infer uniqueness of chain representatives from cohomology
alone.

At either pole the costalk groups are \(E_p,Q,A\) in degrees \(0,1,2\);
their geometric pullback and transfer factors are again \((1,1,n)\) and
\((n,n,1)\). This follows from the actual local support complex
\([V_p\to A\to0 A]\), the local angular degree \(n\), and (GTR2.4) with
supports. On the punctured sphere the factors in degrees \(0,1\) are
\((1,n)\) and \((n,1)\). These maps commute with the entire CSP
localization sequence: the maps \(Q^2\to Q\),
\((x_+,x_-)\mapsto x_+-x_-\), have transfer factor \(n\) on both sides;
the maps \(A\to A^2\to A\), \(a\mapsto(-a,a)\mapsto c_++c_-\), have
factor one throughout their degree-two terms. The endpoint map has
factor \(n\). Thus every support sign and both endpoint kernels survive
the transfer.

\subsection{GTR4. The inverse coefficient map and the two different
adjoints}\label{gtr4.-the-inverse-coefficient-map-and-the-two-different-adjoints}

The already constructed coefficient action is \[
T_a b(u)=b(u/a),
\] \[
\rho_+(a)(f,c_0,c_1)=(f(\cdot/a),c_0,ac_1),\quad
\rho_-(a)(h,d_0,d_1)=(a h(a\,\cdot),ad_0,d_1),
\quad a>0.
\tag{GTR4.1}
\] Direct substitution proves that their inverses are the same displayed
formulas at \(a^{-1}\); they commute with the restrictions as
\(r_p\rho_p(a)=T_ar_p\). This inverse is on coefficient spaces, not on
the geometric cover. It is also the required coefficient map on the
dual: evaluation satisfies \(\lambda(\rho(a)^{-1}\rho(a)x)=\lambda(x)\).
Separation of evaluation proves uniqueness of that coefficient dual map
if evaluation is to be preserved.

Set \(\rho(a)=\rho_+(a)\oplus\rho_-(a)\). The full pullback from CSP is
\[
\mathsf B_n=(\rho(n),T_n,nT_n)
\tag{GTR4.2}
\] on degrees \(0,1,2\). Combine the actual geometric transfer with the
proved inverse coefficient action. Its compatible coefficient
representative is \[
\mathsf U_n=(n\rho(n^{-1}),\ nT_{n^{-1}},\ T_{n^{-1}}).
\tag{GTR4.3}
\] The order of these two operations is immaterial, since the trace is
natural in coefficient maps. On each degree of the displayed complex, \[
\boxed{\mathsf U_n\mathsf B_n=\mathsf B_n\mathsf U_n=n\,1.}
\tag{GTR4.4}
\] The scalar \(n\) therefore comes from the full ramified transfer,
rather than being chosen to obtain a desired spectral formula.
Composition of covers and coefficient actions gives
\(\mathsf U_m\mathsf U_n=\mathsf U_{mn}\).

For the continuous dual complex use CDI's actual Hom differential.
Ordinary transposition of a degree-zero cochain map adds no new sign.
There are three specified maps: \[
\begin{array}{c|ccc}
\text{dual degree}&-2&-1&0\\\hline
\mathsf B_n'&nT_n'&T_n'&\rho(n)'\\
(\mathsf B_n^{-1})'&n^{-1}T_{n^{-1}}'&T_{n^{-1}}'&\rho(n^{-1})'\\
\mathsf U_n'&T_{n^{-1}}'&nT_{n^{-1}}'&n\rho(n^{-1})'.
\end{array}
\tag{GTR4.5}
\] The first is adjoint to original pullback; the second is the
contragredient; the third is adjoint to the constructed transfer with
inverse coefficients. Thus \[
\mathsf U_n'=n(\mathsf B_n^{-1})',\quad
\mathsf B_n'\mathsf U_n'=\mathsf U_n'\mathsf B_n'=n\,1,
\] \[
\langle\mathsf U_n'\lambda,\mathsf B_nx\rangle
=n\langle\lambda,x\rangle.
\tag{GTR4.6}
\] All identities hold on the complete coefficient complex and hence on
its actual cohomology. They do not require a spectral decomposition or a
density assertion. The induced degree-zero dual endpoint characters of
\(\mathsf U_n'\) are \(n,1,1,n\), in the retained order
\(c_0,c_1,d_0,d_1\).

\subsection{GTR5. The original-zeta residue map receives this exact
transfer
action}\label{gtr5.-the-original-zeta-residue-map-receives-this-exact-transfer-action}

Retain the original Mellin comparison \[
\Theta b(s)=\frac12\int_0^\infty b(u)u^s\,\frac{du}{u},\qquad
\Theta T_n b(s)=n^s\Theta b(s),\qquad \Theta J=\mathcal I.
\tag{GTR5.1}
\] Here \(\mathcal B\) is the entire functions rapidly decreasing on
every bounded real strip, \(\mathcal I\) consists of all functions
vanishing to the actual full multiplicity at every nontrivial zero of
original \(\zeta\), and \(\mathcal Q=\mathcal B/\mathcal I\). Write
\(\kappa:Q\to\mathcal Q\), \([b]\mapsto[\Theta b]\). For \(g=[G]\) with
finitely supported full zero jets, RD proves \[
\iota(g)([F])=\sum_{\rho}\operatorname{Res}_{s=\rho}
\frac{F(s)G(1-s)}{\zeta(s)}\,ds.
\tag{GTR5.2}
\] Each member of this finite-jet subspace has a representative
satisfying the required vanishing at every other zero; it is not a
cutoff of the underlying arithmetic. RD constructs those
representatives, proves that the sum is finite, and proves invertibility
of every local multiplicity matrix. Its map \(\iota\) is a linear
isomorphism onto the finite-jet dual inside the entire continuous dual.

The constructed transfer on degree one is \(nT_{n^{-1}}\); hence its
transpose on \(H^{-1}(D^\vee)=Q'\) is \(nT_{n^{-1}}'\). Substituting in
the complete residue integrand gives \[
\boxed{\kappa'\iota\,T_n
=\mathsf U_n'\,\kappa'\iota.}
\tag{GTR5.3}
\] Indeed replacing \(G(1-s)\) by \(n^{1-s}G(1-s)\) is exactly
multiplying the first argument by \(n^{-s}\) and the result by \(n\).
This proves the diagram with the actual geometric transfer. The factor
in RD5 is now received by a constructed sheaf operation, rather than
merely compared as an identical scalar.

For Taylor-coefficient functionals
\(\delta_{\rho,j}([F])=F^{(j)}(\rho)/j!\), the whole formula is \[
\mathsf U_n'\kappa'\delta_{\rho,j}
=n^{1-\rho}\sum_{h=0}^j
\frac{(-\log n)^{j-h}}{(j-h)!}\kappa'\delta_{\rho,h}.
\tag{GTR5.4}
\] It follows by Taylor multiplication, retaining every nilpotent entry
and factorial. No zero simplicity is assumed. The finite-jet image is
invariant, and \(\mathsf U_n'\) is continuous for the weak topology:
evaluation at \(x\) after transposition is evaluation at the already
constructed \(\mathsf U_nx\). Consequently RD's weakly dense receiver
and its actual ambient continuous dual carry compatible transfer
actions. This does not replace that ambient dual by its completion.

The geometric construction applies to every recovered \(n\ge1\). The
corresponding formulas for arbitrary \(a>0\) define the already
available continuous coefficient representation; no nonintegral
holomorphic cover is asserted.

\subsection{GTR6. Normal direction, winding orientation and the retained
extra
factors}\label{gtr6.-normal-direction-winding-orientation-and-the-retained-extra-factors}

The primal normal action on \(H^2\) is \(nT_n\) on \(A\). Write
\(T_a^A\) and \(T_a^Q\) when distinguishing the coefficient space from
its quotient, and retain \(q:A\to Q\). Since \(qT_a^A=T_a^Qq\),
transposition gives \((T_a^A)'q'=q'(T_a^Q)'\). The specified three maps
are \[
\mathsf S_{\mathrm{res}}=n(T_{n^{-1}}^Q)'\quad\text{on }Q',\qquad
\mathsf N_{\mathrm{contra}}=n^{-1}(T_{n^{-1}}^A)'\quad\text{on }A',\qquad
\mathsf N_{\mathrm{transfer}}=(T_{n^{-1}}^A)'\quad\text{on }A'.
\] Their exact comparison, with its domain and receiving inclusion
retained, is \[
q'\mathsf S_{\mathrm{res}}
=n^2\mathsf N_{\mathrm{contra}}q'
=n\mathsf N_{\mathrm{transfer}}q':Q'\longrightarrow A'.
\tag{GTR6.0}
\] Here \(q'\) identifies \(Q'\) with \(J^\perp\); SDT5 proves this is
also a closed strong topological embedding. Thus the residue action has
the displayed extra factors on this specific subspace of the full normal
dual, not on every functional in \(A'\). None of the factors or the
inclusion may be removed.

CDW's winding coordinate \(a[\vartheta]\) maps to \(-a\) in positive
sphere integration. Under that comparison, the GSL normal map is exactly
\[
a[\vartheta]\longmapsto
-[(1-E_+(s))(\Theta a)(s-1)]_{\mathcal I\cap\mathcal I_+},
\tag{GTR6.1}
\] with kernel \(J[\vartheta]\). The sign commutes with both \(nT_n\)
and \(T_{n^{-1}}\), so the trace and pullback tables are unchanged by
this explicit coordinate comparison. It is not permission to omit the
sign from the normal map itself.

The two covers' pole trace maps remain separate. Likewise the
holomorphic swap \(-1/z\) and antiholomorphic swap \(-1/\overline z\)
retain their original CDW angular signs \(-1,+1\) and their CSP
degree-two signs \(+1,-1\). Their current pushforwards are adjoints to
their specified pullbacks. No positive degree is assigned to the
antiholomorphic orientation reversal by (GTR3.2), which concerns only
the holomorphic covers \(b_n\).

\subsection{GTR7. Naturality with the actual lifting
boundary}\label{gtr7.-naturality-with-the-actual-lifting-boundary}

The exact source sequence is \[
0\longrightarrow\mathscr F_J\longrightarrow\mathscr F
\longrightarrow j_!\underline Q\longrightarrow0.
\tag{GTR7.1}
\] The weighted trace is defined on all three terms by the same fibre
multiplicities and is natural in their coefficient maps by GTR2.
Therefore its induced long exact sequences commute with transfer, as
they do with pullback. For the actual connecting map \[
\delta:Q\longrightarrow J(-1),
\tag{GTR7.2}
\] this gives both \[
\delta T_n=nT_n\delta,\qquad
\delta(nT_{n^{-1}})=T_{n^{-1}}\delta.
\tag{GTR7.3}
\] The second relation also follows algebraically from the first because
the specified coefficient actions are invertible. Thus the new trace
construction does not falsely add an independent restriction to the
lifting problem. CW4 and CW7 prove \(\delta=0\) on the entire source;
the present construction proves how this already established lift is
carried through the actual geometric trace and its continuous dual.

The identity (GTR5.4) pairs actual exponents \(\rho\) and \(1-\rho\)
through the full functional equation. It retains their sum and every
multiplicity coordinate. It does not supply an absolute-value estimate
for either eigenvalue. Deriving the support-dual weight bounds needed in
Deligne §3.6 remains a mathematical task; these constructed maps are the
receiving operations for that task.

\subsection{GTR8. Original arithmetic contributions
retained}\label{gtr8.-original-arithmetic-contributions-retained}

The source arithmetic remains \[
\zeta(s)=1+\sum_{k\ge2}k^{-s}
=\prod_p(1-p^{-s})^{-1},\qquad
-\frac{\zeta'}{\zeta}(s)=\sum_p\sum_{m\ge1}(\log p)p^{-ms}
\quad(\Re s>1).
\tag{GTR8.1}
\] The original \(\Sigma\) still has both rational signs and the
complete Mellin identity \[
\mathcal M_0\Sigma f(s)=2\zeta(s)\int_0^\infty f(x)x^s\frac{dx}{x},
\qquad \mathcal M_0=2\Theta.
\tag{GTR8.2}
\] For the source used to construct the global residue isolators, retain
\[
F_0(s)=\frac{s(s-1)}8\pi^{-s/2}\Gamma(s/2)\zeta(s),
\quad F_0(0)=F_0(1)=\frac18,
\] \[
F_0(-2r)=\frac{r(2r+1)(-1)^r\pi^r}{2\,r!}\zeta'(-2r)
=\frac{(1+2r)(2r)}8\pi^{-(1+2r)/2}
\Gamma((1+2r)/2)\zeta(1+2r),\quad r\ge1.
\tag{GTR8.3}
\] RD and SSI give their full unit derivatives, inverse Mellin formula
and trivial-zero Taylor observations. The residue denominator in
(GTR5.2) is original \(\zeta\), not \(F_0\). Its reflection retains \[
\chi(s)=2^s\pi^{s-1}\sin(\pi s/2)\Gamma(1-s)
=\pi^{s-1/2}\frac{\Gamma((1-s)/2)}{\Gamma(s/2)},
\quad\zeta(s)=\chi(s)\zeta(1-s),
\tag{GTR8.4}
\] and the residue reversal has the sign from \(ds=-d(1-s)\), with every
local derivative of \(\chi\) retained as in RD6. The new transfer does
not modify any of these functions or their exceptional contributions.

\subsection{GTR9. Exact result received by the support-duality
programme}\label{gtr9.-exact-result-received-by-the-support-duality-programme}

The factor \(n\) in the original-zeta residue covariance now has the
specific geometric source (GTR2.1): the weighted trace for the actual
ramified cover of the coefficient sheaf. Its inverse coefficient map is
explicit, its two ramified stalks have their full multiplicities, and
its action is calculated on every global, supported and annular group.
The original-zeta finite-jet receiver intertwines its transpose exactly
by (GTR5.3). The adjoint of pullback, the contragredient, and the
adjoint of transfer are linked by (GTR4.6), rather than identified
without their degree factors. This closes that action-comparison step.
It does not identify the ordinary algebraic Verdier dual with the
continuous dual or establish the missing Deligne weight bounds.

\subsection{GTR10. The completed residue receiver and the full trace
action}\label{gtr10.-the-completed-residue-receiver-and-the-full-trace-action}

The completed strong-topology proof SDT0--SDT9 was read in full. Its
accepted source is
\nolinkurl{../tau_weight_cohomology_20260924/CC_STRONG_DUAL_TOPOLOGY_AND_JETS.md},
SHA256
\nolinkurl{7f1a03a6225f1dfc22b167713c0ddbb154635b06f9479f6137b229845b67aefc}.
SDT4 constructs compact lifts through \(A\to Q\); SDT5 proves the strong
topological inclusion \(Q'_\beta\simeq J^\perp\subset A'_\beta\);
SDT6--SDT8 prove strong reflexivity, strong density of the full
finite-jet functionals, and their exact Hausdorff completion. These
results concern the actual coefficient spaces, not bounded samples of
zeros.

Define the residue seminorms on the already constructed finite-support
class space by \[
p_B^{\mathrm{res}}(G)=\sup_{F\in B}\left|
\sum_\rho\operatorname{Res}_{s=\rho}
\frac{F(s)G(1-s)}{\zeta(s)}\,ds\right|,
\qquad B\subset\mathcal Q\text{ bounded}.
\tag{GTR10.1}
\] The sum for each such \(G\) is finite and contains every multiplicity
coordinate in its support. Write \[
\mathcal H_{\mathrm{res}}
=\widehat{(\mathcal Q_{\mathrm{fin}},(p_B^{\mathrm{res}})_B)}.
\tag{GTR10.2}
\] This names the specified Hausdorff completion, not a new topology
silently put on the entire original quotient. SDT8.12 supplies the
topological isomorphism \[
\widehat\iota:\mathcal H_{\mathrm{res}}\xrightarrow{\sim}\mathcal Q'_\beta,
\qquad
\mathfrak I=\kappa'\widehat\iota:
\mathcal H_{\mathrm{res}}\xrightarrow{\sim}Q'_\beta
=H^{-1}(D^\vee).
\tag{GTR10.3}
\] The factor \(1/2\) is still in \(\kappa[a]=[\Theta a]\). To verify
that its transpose is a strong topological isomorphism, the continuous
maps \(\kappa\) and \(\kappa^{-1}\) carry bounded sets to bounded sets.
Each defining strong seminorm therefore pulls back to the seminorm on
the corresponding bounded image. This proves both directions of the
topological assertion in (GTR10.3).

For every \(a>0\), put \(\mathsf V_a=aT_{a^{-1}}\) on
\(\mathcal Q_{\mathrm{fin}}\). Direct substitution into the full
residue, before any completion, proves the exact seminorm identities \[
p_B^{\mathrm{res}}(T_aG)
=a\,p_{T_{a^{-1}}B}^{\mathrm{res}}(G),\qquad
p_B^{\mathrm{res}}(\mathsf V_aG)
=p_{T_aB}^{\mathrm{res}}(G).
\tag{GTR10.4}
\] Indeed the respective multipliers on the second argument at \(1-s\)
are \(a^{1-s}\) and \(a\,a^{-(1-s)}=a^s\). The sets on the right are
bounded because \(T_a\) and its inverse are continuous on
\(\mathcal Q\). Thus both maps and their inverses extend uniquely to
continuous automorphisms \(\widehat T_a\) and \(\widehat{\mathsf V}_a\)
of (GTR10.2). Their products extend those on the dense domain, so \[
\widehat T_a\widehat T_b=\widehat T_{ab},\qquad
\widehat{\mathsf V}_a\widehat{\mathsf V}_b
=\widehat{\mathsf V}_{ab},\qquad
\widehat T_a\widehat{\mathsf V}_a=a\,1.
\tag{GTR10.5}
\] For the recovered integer covers this gives, on the entire
degree-minus-one strong cohomology group, \[
\boxed{
\mathfrak I\widehat T_n=\mathsf U_n'\mathfrak I,
\qquad
\mathfrak I\widehat{\mathsf V}_n=\mathsf B_n'\mathfrak I.}
\tag{GTR10.6}
\] Here \(\mathsf U_n'=nT_{n^{-1}}'\) and \(\mathsf B_n'=T_n'\) in
degree minus one, exactly as in (GTR4.5). The equalities follow on the
dense subspace from (GTR5.3) and (GTR10.4). All maps are continuous into
the Hausdorff space \(Q'_\beta\), so their equalizers are closed; the
equalities hold on the whole completion. This argument proves the full
action without assuming a convergent infinite sum of residues.

The complete pairing is the actual evaluation \[
\widehat{\mathcal R}(F,h)=\widehat\iota(h)(F),
\qquad F\in\mathcal Q,\ h\in\mathcal H_{\mathrm{res}}.
\tag{GTR10.7}
\] For a net \(G_\alpha\to h\) in the specified completion, its original
finite residue pairings converge to (GTR10.7) uniformly on every bounded
\(B\subset\mathcal Q\). This is precisely convergence in the strong dual
under (GTR10.3), so the value is independent of the representing net.
Nondegeneracy in \(h\) follows from injectivity of \(\widehat\iota\).
Nondegeneracy in \(F\) follows from the full zero-jet functionals, which
separate \(\mathcal Q\) by the definition of \(\mathcal I\). No positive
Hermitian form is inferred from this bilinear nondegeneracy.

The full transfer identities become \[
\widehat{\mathcal R}(T_nF,\widehat T_nh)
=n\widehat{\mathcal R}(F,h),\qquad
\widehat{\mathcal R}(F,\widehat{\mathsf V}_nh)
=\widehat{\mathcal R}(T_nF,h).
\tag{GTR10.8}
\] They follow by evaluation of (GTR10.6), or by passing the exact
finite residue identities through the above uniform convergence. The
first retains the geometric degree; the second identifies the transpose
of pullback. All original nilpotent entries of (GTR5.4) persist in this
extension.

Finally, the actual two-support restriction in degree minus one is \[
Q'_\beta\longrightarrow Q'_\beta\oplus Q'_\beta,
\qquad\lambda\longmapsto(\lambda,-\lambda).
\tag{GTR10.9}
\] It is continuous: the seminorm of its value on a bounded pair is
bounded by the sum of the seminorms on its two bounded projections. In
the completed residue coordinates it is exactly \[
h\longmapsto(h,-h).
\tag{GTR10.10}
\] This follows from (GTR10.3) and linearity, retains both separate
poles, and commutes with both actions in (GTR10.6). The embedding into
the overlap dual is explicitly \[
h\longmapsto q'\mathfrak I(h)\in J^\perp\subset A'_\beta,
\quad
(q'\mathfrak I(h))(a)=\widehat{\mathcal R}([\Theta a],h).
\tag{GTR10.11}
\] Its vanishing on \(J\) follows because \(\Theta J=\mathcal I\); SDT5
proves it is a strong topological isomorphism onto that annihilator.
Thus the completed residue receiver maps into the entire signed
support-dual diagram by the original quotient and Mellin maps. Neither
the endpoint groups in degrees zero nor the full normal group in degree
minus two is deleted. This proves the completed trace comparison; it
does not assert that \(\mathcal H_{\mathrm{res}}\) is the original
primal quotient with its original topology, or that its evaluation
pairing is positive.

\subsection{GTR11. Identification of the actual trace on the derived
direct
image}\label{gtr11.-identification-of-the-actual-trace-on-the-derived-direct-image}

There is a stronger receiving statement than the cohomology calculation
of GTR3. Work in the ordinary derived category of complex vector-space
sheaves on the actual three-point space \(X\). CDW4 constructs the
specific quasi-isomorphism \[
\Phi:\Omega\oplus\mathscr W[-1]\xrightarrow{\sim}R\pi_*\mathscr F,
\qquad \mathscr W=j_{\eta!}\underline A_{\{\eta\}},
\tag{GTR11.1}
\] represented by constants and the positive angular form. Here
\(\mathscr W\) is a sheaf on \(X\), not the previously used vector space
\(W=V_+\oplus V_-\). Its only nonzero stalk is \(A\) at \(\eta\). Every
map in the stated representative is continuous on the original
coefficient spaces. CDW5 proves, before taking global cohomology, that
geometric pullback \(P_n\) obeys \[
P_n\Phi=\Phi\begin{pmatrix}1_{\Omega}&0\\0&n1_{\mathscr W[-1]}\end{pmatrix}.
\tag{GTR11.2}
\] Consequently \(P_n\) is an automorphism of this actual derived direct
image. Its inverse is represented through \(\Phi\) by
\(\operatorname{diag}(1,n^{-1})\); no inverse geometric root map is
used.

Let \(\mathsf t_n\) be the morphism on \(R\pi_*\mathscr F\) induced by
the actual sheaf trace of GTR2. Since \(\pi b_n=\pi\), the intermediate
object \(R\pi_*b_{n*}\mathscr F\) is canonically \(R\pi_*\mathscr F\).
To justify this derived identification, \(b_{n*}\) is exact by GTR2 and
preserves injectives because it is right adjoint to the exact
inverse-image functor. An injective resolution therefore computes both
sides of the composite direct image. The sheaf equality (GTR2.4) yields
the equality of derived morphisms \[
\mathsf t_nP_n=n\,1_{R\pi_*\mathscr F}.
\tag{GTR11.3}
\] Explicitly, for the preceding identification
\(\alpha:R\pi_*b_{n*}\mathscr F\xrightarrow{\sim}R\pi_*\mathscr F\), the
two maps are \(P_n=\alpha\,R\pi_*(u_n)\) and
\(\mathsf t_n=R\pi_*(\operatorname{Tr}_n)\alpha^{-1}\). CDW's chain
pullback extends the sheaf unit through its augmentation; it therefore
represents this same \(P_n\). Compose on the right with the already
constructed inverse of \(P_n\). This proves \[
\boxed{
\mathsf t_n=nP_n^{-1},\qquad
\mathsf t_n\Phi
=\Phi\begin{pmatrix}n1_{\Omega}&0\\0&1_{\mathscr W[-1]}\end{pmatrix}.}
\tag{GTR11.4}
\] This identifies the actual trace morphism in the derived category,
including its possible extension information. It does not infer equality
of morphisms merely from equal induced cohomology maps.

The coefficient action of (GTR4.1) commutes with \(\Phi\), preserves
both summands, and has an explicit inverse. Thus the combined pullback
and inverse-coefficient transfer are represented on (GTR11.1) by \[
\mathsf B_n=
\begin{pmatrix}\rho_{\Omega}(n)&0\\0&nT_n\end{pmatrix},
\qquad
\mathsf U_n=
\begin{pmatrix}n\rho_{\Omega}(n^{-1})&0\\0&T_{n^{-1}}\end{pmatrix}.
\tag{GTR11.5}
\] The sheaf morphism \(\rho_{\Omega}(a)\) means
\(\rho_+(a),\rho_-(a),T_a\) on the two chart stalks and the generic
stalk, with the original restrictions retained. Taking derived global
sections and the simultaneous contractions of CW2 gives exactly the full
representatives (GTR4.2)--(GTR4.3). In particular the degree-one factor
\(n\) comes from the \(\Omega\) summand's trace, while the degree-two
factor one comes from its separate angular summand. CDW9's minus between
the angular coordinate and positive sphere integration remains in that
comparison.

Every step is natural in the diagrams for \(J\to A\to Q\), since CDW8's
\(\Phi\) is natural and the trace was constructed naturally in GTR2.
Thus the same derived identification applies simultaneously to the
actual source sequence (GTR7.1), with its boundary maps, rather than to
separately chosen cohomology lists. This strengthens the precise scope
of GTR3 to an actual ordinary derived-sheaf comparison with an explicit
continuous coefficient representative. It makes no assertion of a strict
current pullback through a ramified root map or of uniqueness in an
unspecified category of topological complexes.
